%% kapa: exam=1 round=33 year=2022 period=2 form=A subject=accounting subject_ko=회계학 seq=23 number=63
%% source: Q-Net 감정평가사 기출문제 (cst00309) artlSeq=5209997 file=raw/1008/5209997/2022년도 제33회 감정평가사 1차시험 2교시(감정평가 관계법규, 회계학).pdf
%% conversion: kordoc 4.24.51 -> tools/gen1.py (auto); review=proofread-sonnet; flags=none
\begin{question}{63}
(주)감평은 20×1년 초 부여일로부터 3년의 용역제공을 조건으로 직원 50명에게 각각 주식선택권 10개를 부여하였다. 부여일 현재 주식선택권의 단위당 공정가치는 ￦1,000으로 추정되었으며, 매년 말 추정한 주식선택권의 공정가치는 다음과 같다.
\par{\small\centering\begin{tabular}{|c|c|c|c|}\hline
20×1.12.31 & 20×2.12.31 & 20×3.12.31 & 20×4.12.31 \\ \hline
￦1,000 & ￦1,100 & ￦1,200 & ￦1,300 \\ \hline
\end{tabular}\par}
주식선택권 1개당 1주의 주식을 부여받을 수 있으며 권리가득일로부터 3년간 행사가 가능하다. (주)감평은 20×1년 말과 20×2년 말에 가득기간 중 직원의 퇴사율을 각각 25\%와 28\%로 추정하였으며, 20×1년도와 20×2년도에 실제로 퇴사한 직원은 각각 10명과 2명이다. 20×3년 말 주식선택권을 가득한 직원은 총 35명이다. 20×4년 1월 1일 주식선택권을 가득한 종업원 중 60\%가 본인의 주식선택권 전량을 행사하였을 경우 이로 인한 (주)감평의 자본 증가액은? (단, (주)감평 주식의 주당 액면금액은 ￦5,000이고 주식선택권의 개당 행사가격은 ￦6,000이다.)
\begin{choices}
  \item ￦210,000
  \item ￦420,000
  \item ￦1,050,000
  \item ￦1,260,000
  \item ￦1,470,000
\end{choices}
\end{question}
\answer{4}
